\chapter{SPARQL}
\label{chap:sparql}
Im letzten Kapitel wurde aufgezeigt, wie Inferenz in einer Ontologie mittels der Regelsprache SWRL vollumfänglich genutzt werden kann.

Wie können Informationen aus der Ontologie und den darauf basierenden Inferenzen abgefragt werden?\\
Mit SPARQL --- einer Abfragesrprache --- ist es möglich.

Im folgenden Kapitel, welches, sofern nicht anders vermerkt, auf~\cite{w3sparql_querylang} basiert, stellen wir diese Abfragesprache vor.

Bei SPARQL handelt es sich sich um eine Abfragesprache für RDF.\@ Sie erlaubt es, Abfragen in mehreren Datenquellen vorzunehmen. Dabei werden Anfragen über Graphen vollzogen, auch entlang derer Konjunktionen und Disjunktionen. SPARQL unterstützt weiter Aggregation, Unterabfragen, Negation sowie die Nutzung von Ausdrücken als Werte.

Resultate sind entweder eine Menge von Ergebnissen oder RDF-Graphen.

% (vgl.~\cite[Abstract]{w3sparql_querylang})

\noindent\rule[1ex]{\textwidth}{1pt}
\begin{wrapfigure}[7]{l}{0.1\textwidth}
    \vspace{-12pt}
    \includegraphics[width=0.1\textwidth]{bilder/owl.png}
\end{wrapfigure}
Kennt man die Abfragesprache SQL, so erscheint SPARQL auf den ersten Blick recht ähnlich --- der Name lässt dies bereits vermuten. Dies ist nur teilweise der Fall. Klauseln wie \textit{SELECT} und \textit{WHERE} sind ähnlich. In SPARQL gibt es jedoch keine \textit{FROM}-Klausel.\\
Der Hauptunterschied zwischen SPARQL und SQL ist jedoch die Verwendung von Variablen.  Diese werden in SPARQL in der Regel  bei jeder Abfrage verwendet. Abfragen ohne Variablen kommen in der Praxis (gemäss unserer Erfahrung) eher selten vor.\\
SQL enthält in der Regel sehr wenig Variablen.

\noindent\rule[1ex]{\textwidth}{1pt}

Eine SPARQL-Abfrage besteht aus folgenden Komponenten:
\begin{itemize}
    \item Einem oder mehreren Namespaces
    \item Variablen
    \item Einem (Teil-) Graphen
    \item Einer Menge von Gruppierungen und Aggregationen
    \item Einer Menge von Modifikatoren
    \item Abfragearten
    \item Ausdrücken und Wertevergleichen
\end{itemize}
% (vgl.~\cite[18.1.10 SPARQL Query]{w3sparql_querylang})

\noindent\rule[1ex]{\textwidth}{1pt}
\begin{wrapfigure}[2]{l}{0.1\textwidth}
    \vspace{-19pt}
    \includegraphics[width=0.1\textwidth]{bilder/owl.png}
\end{wrapfigure}
Für die Praxis besteht eine SPARQL-Abfrage in der Regel aus \textit{Namespaces}, einer \textit{SELECT}-Klausel, einer \textit{WHERE}-Klausel sowie \textit{Modifikatoren}.\\

\noindent\rule[1ex]{\textwidth}{1pt}

\section{Beispiel einer SPARQL Abfrage}
\label{sec:sparql_beispiel}

Gegeben sei die folgende Datenbasis:
\lstset{caption={Einfache Datenbasis direkt in SPARQL\protect\footnotemark},captionpos=b}
\begin{lstlisting}
   <http://example.org/book/book1> <http://purl.org/dc/elements/1.1/title> "SPARQL Tutorial".
\end{lstlisting}
\footnotetext{\cite[2.1 Writing a Simple Query]{w3sparql_querylang}}

Die Datenbasis beinhaltet das Objekt \textit{book1} mit dem Attribut \textit{title}, welches den Wert \textit{``SPARQL Tutorial''} enthält.

Möchte ich nun den Titel des Buches abfragen, so kann ich folgende Abfrage verwenden:
\lstset{caption={Beispiel einer einfachen SPARQL-Abfrage\protect\footnotemark},captionpos=b}
\begin{lstlisting}
    SELECT
        ?title
    WHERE {
        <http://example.org/book/book1> <http://purl.org/dc/elements/1.1/title> ?title .
    }
\end{lstlisting}
\footnotetext{\cite[2.1 Writing a Simple Query]{w3sparql_querylang}}

Sie ergibt folgendes Resultat:
\noindent\hspace*{15mm}
\begin{table}[h]
    \centering
    \begin{tabular}{|l|}
        \hline
        \multicolumn{1}{|c|}{\textbf{title}} \\ \hline
        ``SPARQL Tutorial''                    \\ \hline
    \end{tabular}
    \caption{Resultat einer einfachen SPARQL-Abfrage\protect\footnotemark}
\end{table}
\footnotetext{\cite[2.1 Writing a Simple Query]{w3sparql_querylang}}

\section{Namespaces}
\label{sec:sparql_namespaces}

SPARQL ist eine Abfragesprach für RDF, das seinerseits auf XML basiert. Dadurch sind XML-Namespaces auch in SPARQL verfügbar. Bei Namespaces handelt es sich um Referenzen auf andere XML-basierte Dokumente, welche modular genutzt werden können.\\
Ein Namespace besteht aus einem Kürzel (welches auch leer sein kann), sowie aus der kompletten Adresse der Ressource (vgl.~\cite[2.1 Introduction]{w3rdf_syntax}).

\lstset{language=XML,
  caption={Beispiel von Namespaces in RDF\protect\footnotemark},
  frame=L,
  basicstyle=\small\normalfont\sffamily,  % the size of the fonts that are used for the code
  stepnumber=1,                           % the step between two line-numbers.
                                          % If it is 1 each line will be numbered
  numbersep=10pt,                         % how far the line-numbers are from the code
  tabsize=2,                              % tab size in blank spaces
  extendedchars=true,                     %
  breaklines=true,                        % sets automatic line breaking
  captionpos=b,                           % sets the caption-position to bottom
  mathescape=true,
  showspaces=false,                       % Leerzeichen anzeigen ?
  showtabs=false,                         % Tabs anzeigen ?
  xleftmargin=17pt,
  framexleftmargin=17pt,
  framexrightmargin=17pt,
  framexbottommargin=5pt,
  framextopmargin=5pt,
  showstringspaces=false                  % Leerzeichen in Strings anzeigen ?
  belowcaptionskip=5em,
  belowskip=3em,
  aboveskip=3em
 }
\begin{lstlisting}
$\textless$rdf:RDF xmlns:rdf="http://www.w3.org/1999/02/22-rdf-syntax-ns#"
         xmlns:dc="http://purl.org/dc/elements/1.1/"
         xmlns:ex="http://example.org/stuff/1.0/"$\textgreater$
\end{lstlisting}
\footnotetext{\cite[2.6 Completing the Document: Document Element and XML Declaration]{w3rdf_syntax}}

\section{Variablen}
\label{sec:sparql_variablen}
Bei Abfragen unterstützt SPARQL die Verwendung von Variablen. Diese können dabei an einer beliebigen Stelle, wie bei (Teil-) Graphen, Gruppierungen, Aggregationen und Modifikatoren eingesetzt werden.\\
Die Variablen dienen als Platzhalter für abzufragende Objekte, Subjekte und Prädikate. Variablen beginnen entweder mit einem Dollarzeichen \textit{\$} oder einem Fragezeichen \textit{?}, welche selbst nicht Teil der Variable sind.

\lstset{caption={Beispiel einer SPARQL Abfrage mit den Variablen \textit{?object}, \textit{?predicate} und \textit{?subject}.},captionpos=b}
\begin{lstlisting}
    SELECT DISTINCT
        *
    WHERE {
      ?object ?predicate ?subject
    }
    LIMIT
        10
\end{lstlisting}

Wird obige Anfrage auf einen SPARQL-Endpunkt, wie z.B. \textit{http://dbpedia.org/}, angewendet, so ergibt dies folgendes Resultat:
\noindent\hspace*{15mm}
\begin{table}[h]
    \centering
    \begin{tabular}{l|l|l|}
        \hline
        \multicolumn{1}{|l|}{\textbf{object}}                                & \textbf{predicate}                               & \textbf{subject}                                  \\ \hline
        \multicolumn{1}{|l|}{http://dbpedia.org/ontology/acceleration}      & ... rdf-syntax-ns\#type & ... owl\#FunctionalProperty \\ \hline
        \multicolumn{1}{|l|}{http://dbpedia.org/ontology/torqueOutput}      & ... rdf-syntax-ns\#type & ... owl\#FunctionalProperty \\ \hline
        \multicolumn{1}{|l|}{http://dbpedia.org/ontology/birthDate}         & ... rdf-syntax-ns\#type & ... owl\#FunctionalProperty \\ \hline
        \multicolumn{1}{|l|}{http://dbpedia.org/ontology/birthYear}         & ... rdf-syntax-ns\#type & ... owl\#FunctionalProperty \\ \hline
        \multicolumn{1}{|l|}{http://dbpedia.org/ontology/deathDate}         & ... rdf-syntax-ns\#type & ... owl\#FunctionalProperty \\ \hline
        \multicolumn{1}{|l|}{http://dbpedia.org/ontology/deathYear}         & ... rdf-syntax-ns\#type & ... owl\#FunctionalProperty \\ \hline
        \multicolumn{1}{|l|}{http://dbpedia.org/ontology/diameter}          & ... rdf-syntax-ns\#type & ... owl\#FunctionalProperty \\ \hline
        \multicolumn{1}{|l|}{http://dbpedia.org/ontology/displacement}      & ... rdf-syntax-ns\#type & ... owl\#FunctionalProperty \\ \hline
        \multicolumn{1}{|l|}{http://dbpedia.org/ontology/height}            & ... rdf-syntax-ns\#type & ... owl\#FunctionalProperty \\ \hline
        \multicolumn{1}{|l|}{http://dbpedia.org/ontology/latestReleaseDate} & ... rdf-syntax-ns\#type & ... owl\#FunctionalProperty \\ \hline
    \end{tabular}
    \caption{Resultat einer Abfrage des DBPedia-Endpunktes limitiert auf 10 Ergebnisse.}
\end{table}

\section{(Teil-) Graphen}
\label{sec:sparql_graph}

SPARQL basiert auf dem Matching eines (Teil-) Graphen in einem gegebenen Graphen. Dabei ist ein (Teil-) Graph eine Menge von Tripeln.

Man unterscheidet zwischen den folgenden Arten von Teilgraphen:
\begin{itemize}
\item \textit{Grundlegende Teilgraphen}

Hierbei muss eine Menge von Tripeln derjenigen im gegebenen Graphen entsprechen

\item \textit{Gruppierende Teilgraphen}

Hierbei müssen \textit{alle} Elemente einer Menge von Tripeln von grundlegenden Teilgraphen wiederum derjenigen im gegebenen Graphen entsprechen

\item \textit{Optionale Teilgraphen}

Zusätzliche Teilgraphen können die Lösungsmenge erweitern

\item \textit{Alternative Teilgraphen}

Versuch, ob zwei oder mehrere mögliche Teilgraphen denjenigen im gegebenen Graphen entsprechen

\item \textit{Teilgraphen in namensbasierten Graphen}

Versuch, ob Teilgraphen denjenigen im gegebenen, namensbasierten Graphen entsprechen
\end{itemize}

Ein grundlegender Teilgraph ist somit eine Menge von Tripeln, wobei jeder Tripel gemäss RDF-Spezifikation immer aus Subjekt, Prädikat und Objekt bestehen muss (vgl.~\cite[3.1 Triples]{w3rdf}).

\lstset{caption={Beispiel eines grundlegenden Teilgraphen in SPARQL}}
\begin{lstlisting}
    ?subject ?predicate ?object.
\end{lstlisting}

\noindent\rule[1ex]{\textwidth}{1pt}
\begin{wrapfigure}[3]{l}{0.1\textwidth}
    \vspace{-16pt}
    \includegraphics[width=0.1\textwidth]{bilder/owl.png}
\end{wrapfigure}
Ein Teilgraph wird in SPARQL immer zur Gruppierung einer Abfrage verwendet. Dies entspricht in der SQL-Sprache der \textit{WHERE}-Klausel.\\
Ein Teilgraph wird immer mit dem Punkt-Operator beendet.

\noindent\rule[1ex]{\textwidth}{1pt}

Ein gruppierender Teilgraph wird in der SPARQL-Abfrage immer mit geschweiften Klammern \textit{\{\}} umschlossen. Es ist möglich die Resultate einer Gruppierung zu filtern. Dazu wird \textit{FILTER} als Schlüsselwort verwendet. Dieses bezieht sich immer auf die Gruppierung, in welcher es steht.

\lstset{caption={Beispiel eines gruppierenden Teilgraphen mit dem \textit{FILTER}-Schlüsselwort\protect\footnotemark}}
\begin{lstlisting}
    {
        ?x foaf:name ?name.
        ?x foaf:mbox ?mbox.
        FILTER regex(?name, "Smith")
    }
\end{lstlisting}
\footnotetext{\cite[5.2.2 Scope of Filters]{w3sparql_querylang}}

Im obigem Beispiel werden alle Objekte und die Werte ihrer Attribute (\textit{name} und \textit{mbox}, welche im Namespace \textit{foaf} bezeichnet werden) abgefragt.

Dabei kann eine komplette leere Menge --- wenn keine Objekte vorhanden sind --- oder aber eine leere Teilmenge --- wenn ein Objekt z.B. nicht über die beiden Attribute \textit{name} und \textit{mbox} verfügt --- zurückgegeben werden.

Es bleiben nur Objekte zurück, welche über das Attribut \textit{foaf:name} verfügen und dessen Inhalt \textit{``Smith''} ist.

\section{Gruppierungen und Aggregationen}
\label{sec:sparql_gruppierungenaggregationen}
Aggregationen bezeichnen eine bestimmte Art der Verbindung zwischen Objekten.~\cite{wiki:aggregation} Aggregationen wenden Ausdrücke bzw. Funktionen auf eine Lösungsmenge an. Dabei enthält eine Lösung normalerweise eine einzelne Gruppe mit allen Lösungen.

Gruppierungen werden mittels \textit{GROUP BY} angegeben.

SPARQL unterstützt aktuell die Aggregationen \textit{COUNT}, \textit{SUM}, \textit{MIN}, \textit{MAX}, \textit{AVG}, \textit{GROUP\_CONCAT} und \textit{SAMPLE}.

Die Aggregation wird dann verwendet, wenn ein Resultat über eine Gruppe von Lösungen berechnet werden soll (Beispiel: Maximaler Wert einer bestimmten Variablen).

\lstset{caption={Beispiel einer Abfrage mit Gruppierung und Aggregation\protect\footnotemark}}
\lstset{language=SQL}
\begin{lstlisting}
    PREFIX books: $\textless$http://books.example/$\textgreater$
    SELECT (
        SUM(?lprice) AS ?totalPrice
    )
    WHERE {
      ?org books:affiliates ?auth .
      ?auth books:writesBook ?book .
      ?book books:price ?lprice .
    }
    GROUP BY ?org
    HAVING (SUM(?lprice) $\textgreater$ 10)
\end{lstlisting}
\footnotetext{\cite[11.1 Aggregate Example]{w3sparql_querylang}}
    
\section{Modifikatoren}
\label{sec:sparql_modifikatoren}
Abfragen mittels SPARQL generieren eine ungeordnete Menge von Lösungen. Jede dieser Lösungen stellt eine Funktion von Variablen zu RDF-Termen dar. So gewonnene Lösungen werden als Sequenz von Lösungen behandelt, zunächst ohne spezifische Ordnung.

Um eine Sequenz von Lösungen zu ordnen, werden Modifikatoren verwendet.\\
Dabei unterscheidet man zwischen folgenden Modifikatoren:

\begin{itemize}
    \item \textit{Order}

        Sortiert die Lösungen auf- oder absteigend nach einer bestimmten Variablen

    \item \textit{Projection}

        Erlaubt die Auswahl von spezifischen Variablen mittels der \textit{Select}-Klausel

    \item \textit{Distinct}

        Stellt sicher, dass die Lösungen in der Lösungsmenge einmalig sind

    \item \textit{Reduced}

        Verhindert die Elimination von Duplikaten in der Lösungsmenge

    \item \textit{Offset}

        Definiert, ab welchem Element der Lösungsmenge die gefundenen Lösungen ausgegeben werden

    \item \textit{Limit}

        Definiert die maximale Anzahl an gefundenen Lösungen, die ausgegeben werden soll

\end{itemize}

Die \textit{Modifikatoren} werden in der Reihenfolge der obigen Liste angewendet.


\section{Abfragearten}
\label{sec:sparql_abfragearten}

SPARQL unterscheidet zwischen vier Abfragearten:
\begin{itemize}
    \item \textit{SELECT}

        Liefert alle Variablen oder eine Teilmenge derselben eines Graphen

    \item \textit{CONSTRUCT}

        Liefert einen RDF-Graphen durch Substitution der Variablen in einer Menge von gegebenen Tripeln

    \item \textit{ASK}

        Liefert einen Wahrheitswert (Boolean), welcher angibt, ob die Menge von angefragten Tripeln derjenigen in dem gegebenen Graphen entspricht

    \item \textit{DESCRIBE}

        Liefert einen RDF-Graphen, welcher die gefundenen Ressourcen beschreibt
\end{itemize}

\subsection{SELECT}
\label{subsec:sparql_abfragearten_select}
Die \textit{SELECT}-Abfrageart liefert direkt Variablen und deren Belegung, in Form von Projektionen. Dabei entstehen neue Variablen-Belegungen. Die spezifischen, gewünschten Variablen werden in Form einer Liste, durch Leerzeichen getrennt, angegeben.

\lstset{caption={Beispiel der \textit{SELECT}-Abfrageart in SPARQL},captionpos=b,language=SQL}
\begin{lstlisting}
    SELECT ?a ?b ?c
    WHERE {
        ...
\end{lstlisting}

Die $ SELECT * $-Syntax ist eine Kurzschreibweise und bedeutet die Verwendung bzw.\ die Auflistung aller in der Abfrage verwendeten Variablen. Davon ausgenommen sind Variablen innerhalb von \textit{FILTER}-Anweisungen sowie Variablen, welche rechts des Minus-Operators stehen. Die Syntax darf nur verwendet werden, wenn die Abfrage keine \textit{GROUP BY}-Aggregation verwendet.
\lstset{caption={Beispiel der \textit{SELECT *} Kurzschreibweise in SPARQL},captionpos=b,language=SQL}
\begin{lstlisting}
    SELECT *
    WHERE {
    ...
\end{lstlisting}

Durch die \textit{SELECT}-Abfrageart lassen sich neuen Variablen durch Ausdrücke einführen. Diese werden mittels $(SELECT \; expr \; AS \; var)$ eingeführt.

\lstset{caption={Beispiel eines Ausdrucks der \textit{SELECT}-Abfrageart},captionpos=b,language=SQL}
\begin{lstlisting}
SELECT ?title (?p * (1 - ?discount) AS ?price)
WHERE {
...
\end{lstlisting}


\subsection{CONSTRUCT}
\label{subsec:sparql_abfragearten_construct}
Die \textit{CONSTRUCT}-Abfrageart gibt als Antwort einen einzigen, durch eine Vorlage definierten RDF-Graphen zurück: Jede Variable der Abfrage wird durch die passende Lösung aus der Lösungssequenz ersetzt. Die Lösungsmenge wird schliesslich durch Vereinigung der Tripel auf einen einzigen RDF-Graphen reduziert.

\lstset{caption={Beispiel einer SPARQL-Abfrage unter Nutzung der \textit{CONSTRUCT}-Abfrageart\protect\footnotemark},captionpos=b}
\begin{lstlisting}
PREFIX foaf:    $\textless$http://xmlns.com/foaf/0.1/$\textgreater$
PREFIX vcard:   $\textless$http://www.w3.org/2001/vcard-rdf/3.0#$\textgreater$
CONSTRUCT   { $\textless$http://example.org/person#Alice$\textgreater$ vcard:FN ?name }
WHERE       { ?x foaf:name ?name }
\end{lstlisting}
\footnotetext{\cite[16.2 CONSTRUCT]{w3sparql_querylang}}

Die obige Abfrage erzeugt folgende Ausgabe:
\lstset{caption={Beispiel der Ausgabe einer SPARQL-Abfrage unter Nutzung der \textit{CONSTRUCT}-Abfrageart\protect\footnotemark},captionpos=b}
\begin{lstlisting}
@prefix vcard: $\textless$http://www.w3.org/2001/vcard-rdf/3.0#$\textgreater$ .

$\textless$http://example.org/person#Alice$\textgreater$ vcard:FN "Alice" .
\end{lstlisting}
\footnotetext{\cite[16.2 CONSTRUCT]{w3sparql_querylang}}

Entsteht bei diesem Vorgang ein ungültiges Tripel, wird dieses nicht im Ausgabegraphen angezeigt.\\
Enthält die Vorlage eines Graphen Tripel ohne Variablen, erscheinen diese dennoch im Ausgabegraphen.

\lstset{caption={Beispiel einer SPARQL Abfrage mit Nutzung der \textit{CONSTRUCT}-Abfrageart unter Verwendung leerer Knoten\protect\footnotemark},captionpos=b}
\begin{lstlisting}
PREFIX foaf:    $\textless$http://xmlns.com/foaf/0.1/$\textgreater$
PREFIX vcard:   $\textless$http://www.w3.org/2001/vcard-rdf/3.0#$\textgreater$

CONSTRUCT { ?x  vcard:N _:v .
            _:v vcard:givenName ?gname .
            _:v vcard:familyName ?fname }
WHERE
 {
    { ?x foaf:firstname ?gname } UNION  { ?x foaf:givenname   ?gname } .
    { ?x foaf:surname   ?fname } UNION  { ?x foaf:family_name ?fname } .
 }
\end{lstlisting}
\footnotetext{\cite[16.2.1 Templates with Blank Nodes]{w3sparql_querylang}}

Die obige Abfrage erzeugt folgende Ausgabe:
\lstset{caption={Beispiel der Ausgabe einer SPARQL Abfrage mit Nutzung der \textit{CONSTRUCT}-Abfrageart unter Verwendung leerer Knoten\protect\footnotemark},captionpos=b}
\begin{lstlisting}
@prefix vcard: $\textless$http://www.w3.org/2001/vcard-rdf/3.0#$\textgreater$ .

_:v1 vcard:N         _:x .
_:x vcard:givenName  "Alice" .
_:x vcard:familyName "Hacker" .

_:v2 vcard:N         _:z .
_:z vcard:givenName  "Bob" .
_:z vcard:familyName "Hacker" .
\end{lstlisting}
\footnotetext{\cite[16.2.1 Templates with Blank Nodes]{w3sparql_querylang}}

\subsection{ASK}
\label{subsec:sparql_abfragearten_ask}
Die \textit{ASK}-Abfrageart beantwortet in der Ausgabe, ob eine Lösung für die gegebene Abfrage existiert.

Gegeben sei die folgende Datenbasis:
\lstset{caption={Einfache Datenbasis direkt in SPARQL\protect\footnotemark},captionpos=b}
\lstset{language=XML}
\begin{lstlisting}
@prefix foaf:       $\textless$http://xmlns.com/foaf/0.1/$\textgreater$ .

_:a  foaf:name       "Alice" .
_:a  foaf:homepage   $\textless$http://work.example.org/alice/$\textgreater$ .

_:b  foaf:name       "Bob" .
_:b  foaf:mbox       $\textless$mailto:bob@work.example$\textgreater$ .
\end{lstlisting}
\footnotetext{\cite[16.3 ASK]{w3sparql_querylang}}

Wird nun folgende Abfrage getätigt:
\lstset{caption={Beispiel einer SPARQL-Abfrage mit Nutzung der \textit{ASK}-Abfrageart\protect\footnotemark},captionpos=b}
\lstset{language=XML}
\begin{lstlisting}
    PREFIX foaf:    $\textless$http://xmlns.com/foaf/0.1/$\textgreater$
    ASK  { 
        ?x foaf:name  "Alice"
    }
\end{lstlisting}
\footnotetext{\cite[16.3 ASK]{w3sparql_querylang}}
So antwort das System: $true$, da die Datenbasis ein Individuum mit dem Attribut $foaf:name$ enthält, dessen Wert $``Alice''$ ist.

\subsection{DESCRIBE}
\label{subsec:sparql_abfragearten_describe}
Bei der \textit{DESCRIBE}-Abfrageart antwortet das System mit einem RDF-Graphen, der Informationen in Form von RDF-Daten über Ressourcen enthält. Als Ressourcen kommen entweder \textit{IRI}s oder Variablen in Betracht.

\lstset{caption={Beispiel der \textit{DESCRIBE}-Abfrageart in SPARQL},captionpos=b,language=SQL}
\begin{lstlisting}
    DESCRIBE
        ?a ?b ?c
    WHERE {
    ...
\end{lstlisting}

Die $ DESCRIBE * $-Syntax ist die Kurzschreibweise und bedeutet die Verwendung bzw.\ die Auflistung aller in der Abfrage verwendeten Variablen.
\lstset{caption={Beispiel der \textit{DESCRIBE *} Kurzschreibweise in SPARQL},captionpos=b,language=SQL}
\begin{lstlisting}
DESCRIBE *
WHERE {
...
\end{lstlisting}

\section{Ausdrücke und Wertevergleiche}
\label{sec:sparql_ausdruecke}
In SPARQL kann jede Abfrage in SPARQL mit einem Filter versehen werden. Dabei wird versucht ein bestimmtes Muster (durch Beschränkungen) auf Graphen anzuwenden um so die Graphen einzuschränken (filtern).

Jedes RDF-Literal kann durch einen bestimmten Datentyp definiert sein, wie zum Beispiel Wahrheitswerte (Boolean) oder Datum und Uhrzeit (DateTime). Der Filter in SPARQL erlaubt dabei den Wertevergleich auf typisierte RDF-Literale. Die dabei verwendeten Operanden und Datentypen finden sich unter den Webseiten~\href{http://www.w3.org/TR/xpath-functions/}{w3.org/TR/xpath-functions}\footnote{\url{http://www.w3.org/TR/xpath-functions/}} und~\href{http://www.w3.org/TR/sparql11-query/\#operandDataTypes}{w3.org/TR/sparql11-query/}\footnote{\url{http://www.w3.org/TR/sparql11-query/\#operandDataTypes}}.

Eine SPARQL-Abfrage unter Verwendung eines Filters kann wie folgt aussehen:
\lstset{caption={Beispiel einer einfachen SPARQL-Abfrage unter Verwendung eines Filters},captionpos=b}
\begin{lstlisting}
    ...
    SELECT
        ?book, ?amount
    WHERE {
      ?book :isInLibrary :BerneUniversityLibrary
      ?book :hasAmount ?amount
      FILTER (?amount $\textgreater$ xsd:integer(10))
    }
\end{lstlisting}
In diesem Beispiel wird eine Datenbank nach mehrfachen Vorkommen (mindestens 10 Mal) eines Buchexemplares in der Bibliothek der Universität Bern abgefragt. Die Anzahl der Exemplare eines Buches wird durch die Variable $?amount$ zurückgegeben.

Ist man in obigem Beispiel nicht sicher, welchem Datentyp die Variable $?amount$ entspricht, kann man die Funktionalität zur Umwandlung (Casting, Typenumwandlung) von Variablen nützen:
\lstset{caption={Beispiel einer einfachen SPARQL-Abfrage unter Verwendung eines Filters mit Typenumwandlung},captionpos=b}
\begin{lstlisting}
    ...
    SELECT
        ?book, ?amount
    WHERE {
      ?book :isInLibrary :BerneUniversityLibrary
      ?book :hasAmount ?amount
      FILTER (xsd:integer(?amount) $\textgreater$ xsd:integer(10))
    }
\end{lstlisting}

Die genaue Auswertung eines Filters lassen sich unter der Webseite~\href{http://www.w3.org/TR/sparql11-query/\#evaluation}{w3.org/TR/sparql11-query/\#evaluation}\footnote{\url{http://www.w3.org/TR/sparql11-query/\#evaluation}}, die interne Verwendung und Auswertung von Operatoren unter der Webseite~\href{http://www.w3.org/TR/sparql11-query/\#OperatorMapping}{w3.org/TR/sparql11-query/\#OperatorMapping}\footnote{\url{http://www.w3.org/TR/sparql11-query/\#OperatorMapping}} nachlesen.

\subsection{Funktionen}
\label{subsec:sparql_ausdruecke_funktionen}

SPARQL bietet eine Vielzahl an Operatoren und Funktionen. Diese einzeln aufführen und beschreiben zu wollen, würde den Rahmen dieser Arbeit sprengen. Daher werden nur die Operatoren und Funktionen aufgeführt, welche in der vorliegenden Arbeit hauptsächlich angewendet wurden. Eine genauere Beschreibung findet sich unter der Webseite~\href{http://www.w3.org/TR/sparql11-query/\#SparqlOps}{w3.org}\footnote{\url{http://www.w3.org/TR/sparql11-query/\#SparqlOps}}.

\subsubsection{BOUND}
\label{subsec:sparql_ausdruecke_funktionen_bound}
Mittels \textit{BOUND} lässt sich prüfen, ob ein Objekt eine Wert-Bindung hat. Man kann es als eine Art der Überprüfung auf gewisse Eigenschaften bezeichnen.

Möchte man zum Beispiel wissen, welche Personen über eine Relation zu einem Datum (z.B. ihr Erfassungsdatum) haben, so bestimmt man dies mittels:
\lstset{caption={Beispiel einer simplen Ontologie zur Nutzung der \textit{BOUND}-Funktion in einer Abfrage\protect\footnotemark},captionpos=b}
\begin{lstlisting}
@prefix foaf:        $\textless$http://xmlns.com/foaf/0.1/$\textgreater$ .
@prefix dc:          $\textless$http://purl.org/dc/elements/1.1/$\textgreater$ .
@prefix xsd:          $\textless$http://www.w3.org/2001/XMLSchema#$\textgreater$ .

_:a  foaf:givenName  "Alice".

_:b  foaf:givenName  "Bob" .
_:b  dc:date         "2005-04-04T04:04:04Z"^^xsd:dateTime .
\end{lstlisting}
\footnotetext{\cite[17.4.1.1 BOUND]{w3sparql_querylang}}

\lstset{caption={Beispiel einer SPARQL-Abfrage zur Nutzung der \textit{BOUND}-Funktion\protect\footnotemark},captionpos=b}
\begin{lstlisting}
    PREFIX foaf: $\textless$http://xmlns.com/foaf/0.1/$\textgreater$
    PREFIX dc:   $\textless$http://purl.org/dc/elements/1.1/$\textgreater$
    PREFIX xsd:   $\textless$http://www.w3.org/2001/XMLSchema#$\textgreater$
    SELECT
        ?givenName
    WHERE {
        ?x foaf:givenName  ?givenName .
        OPTIONAL { ?x dc:date ?date } .
        FILTER ( bound(?date) )
    }
\end{lstlisting}
\footnotetext{\cite[17.4.1.1 BOUND]{w3sparql_querylang}}

Die obige Abfrage ergibt folgendes Resultat:
\noindent\hspace*{15mm}
\begin{table}[h]
    \centering
    \begin{tabular}{|l|}
        \hline
        \multicolumn{1}{|c|}{\textbf{givenName}} \\ \hline
        ``Bob''                    \\ \hline
    \end{tabular}
    \caption{Resultat einer SPARQL-Abfrage mittels \textit{BOUND}-Filter\protect\footnotemark}
\end{table}
\footnotetext{\cite[17.4.1.1 BOUND]{w3sparql_querylang}}

Es wird nur die Person \textit{``Bob''} gefunden, da nur bei dieser Person eine Datum-Relation definiert ist.

\subsubsection{IF}
\label{ssubsec:sparql_ausdruecke_funktionen_if}
Die \textit{IF}-Funktion gibt Argumente aufgrund einer Bedingung, welche selbst auch ein Argument ist, zurück. Die Funktion verwendet drei Argumente als Parameter. Die Argumente sind:
\begin{itemize}
    \item Argument 1: Bedingung
    \item Argument 2: Rückgabewert wenn die Bedingung (Argument 1) zutrifft
    \item Argument 3: Rückgabewert wenn die Bedingung (Argument 1) nicht zutrifft
\end{itemize}

\lstset{caption={Beispiel der \textit{IF}-Funktion in SPARQL},captionpos=b,language=XML}
\begin{lstlisting}
IF(?a = 2,?b,?c)
\end{lstlisting}

In prozeduraler Programmierung ausgedrückt:
\lstset{caption={Beispiel der \textit{IF}-Funktion in SPARQL anhand prozeduraler Programmierung ausgedrückt},
    captionpos=b,
    language=Java
}
\begin{lstlisting}
if (a == 2) {
    return b;
} else {
    return c;
}
\end{lstlisting}

\subsubsection{IN}
\label{subsec:sparql_ausdruecke_funktionen_in}
Die \textit{IN}-Funktion prüft, ob ein Element in einer Liste vorkommt. Das zu prüfende Element wird der \textit{IN}-Funktion vorangestellt, danach folgt die zu prüfende Liste. Kommt das Element in der Liste vor, gibt die Funktion den Wahrheitswert $wahr$ bzw. $true$ zurück, andernfalls $falsch$ bzw. $false$.

\lstset{caption={Beispiel der \textit{IN}-Funktion in SPARQL},captionpos=b,language=SQL}
\begin{lstlisting}
2 IN (1, 2, 3) % Wahr
2 IN ()        % Falsch
\end{lstlisting}
\footnotetext{\cite[17.4.1.9 IN]{w3sparql_querylang}}
In den obigen Beispielen ist der erste Ausdruck wahr, da 2 in der gegebenen Liste von (1, 2, 3) vorkommt. Der zweite Ausdruck ist falsch, da 2 nicht in der leeren Liste () vorkommt.

Analog zur \textit{IN}-Funktion existiert die \textit{NOT IN}-Funktion. Dabei handelt es sich um die Negation der \textit{IN}-Funktion.


\noindent\rule[1ex]{\textwidth}{1pt}
\begin{wrapfigure}[7]{l}{0.1\textwidth}
    \vspace{-12pt}
    \includegraphics[width=0.1\textwidth]{bilder/elephant.png}
\end{wrapfigure}
Wie beschrieben sind jetzt alle Komponenten vorhanden, um Ausflüge abbilden zu können, die den in~\autoref{chap:graph_data} (\nameref{chap:graph_data}) genannten Kriterien entsprechen.

In diesem Kapitel haben wir die bisher fehlende Abfragesprache abgehandelt.

Um auf das beschriebene Beispiel der Familie, welche einen Ausflug plant, zurückzukommen: Mit der nachfolgenden Abfrage gelingt es, alle Ausflüge mit den Kriterien  \textit{familienfreundlich}, \textit{regional} und \textit{actionreich} zu finden.

\begin{lstlisting}[caption={Beispiel einer Abfrage um alle familiengerechten, actionreichen und regionalen Ausflüge der Ontologie auszugeben},captionpos=b,language=SQL]
    SELECT
        *
    WHERE {
        ?object :familiengerecht true;
            :regional true;
            :action true.
    }
\end{lstlisting}

\noindent\rule[1ex]{\textwidth}{1pt}
